\documentclass[10pt
 ,aspectratio=1610
]{beamer}	%% 16:10 ratio
% ==================================================================
%% load settings

\usepackage{xcolor}
\usepackage{eurosym}
\usepackage{comment}
\newcommand\cropped[1]{%
	\immediate\write18{convert -trim #1.png #1cropped.png}%
	\includegraphics[width=0.8\linewidth]{#1cropped.png}}
% ==================================================================
%% set up title page data
\title[] {Optimierungsmethoden \\[1mm] für nachhaltige Baukonstruktionen \\[1.5mm] Homework assignments }

%\subtitle{Einführung}

\author[M.~Gfrerer,  A.~Borkovic] % (optional)
{Michael Gfrerer, Aleksandar Borkovic}

\date % (optional)
{24. März 2026}
% ==================================================================
%++++++++++++++++++++
% the document
%++++++++++++++++++++
%
\begin{document}

%To start with page number 1 after titleframe
\addtocounter{framenumber}{-1}
\mytitle  % Dieser Befehl anstatt \maketitle ist wichtig, da sonst
          % % nicht das Hintergrundbild gesetzt wird
% ==================================================================
%{\frame{\frametitle{Inhalt}
%    \tableofcontents}}
%\input{content/einfuehrung}
% ---------------------------------------------------------------------------
% ---------------------------------------------------------------------------
\begin{comment}
\begin{frame}[c]
	\frametitle{Exercises}
	\begin{enumerate}
		\item Optimization of frames
		\item Size optimization of trusses
		\item Optimization of building insulation
		\item Size optimization of networks
		\item Worst case loading
		\begin{itemize}
			\item Truss
			\item FEM for plane structures
			\item analytic methods for frames
            \item two-span beam
		\end{itemize}
        \begin{itemize}
            \item dead and live loads
        \end{itemize}
		%\item Optimization of retainment walls
	\end{enumerate}
	
\end{frame}  
\end{comment}
% ---------------------------------------------------------------------------

% ---------------------------------------------------------------------------

% ==================================================================


% ==================================================================


% ==================================================================


% ==================================================================

% 

\begin{frame}[c]
\frametitle{Topic: FEM}
Title: Optimize the position of stiffeners of a plate
\vfill
Tasks:
\begin{itemize}
    \item Learn about the bending of thin plates
    \item Recall the beam bending theory
    \item Learn about approximating a stiffened plate as a bending plate with carefully added stiffness of beams
    \item Introduce yourself to the basics of shape optimization
    \item Introduce yourself to the basics of FreeFEM 
    \item Make yoursef familiar with the given code
    %\item Possible improvement: Enhance the code so that it can address a free vibration problem, and define the lowest eigenfrequency as the objective function 
  %  \item Decision variables are positions of stiffeners. 

\end{itemize}

In the current code, compliance (external work) due to a uniformly distributed load is defined as an objective function.

\end{frame}

\begin{frame}[c]

Deliverables:
\begin{itemize}

    \item Define the analysis cases: for example, consider cases with different numbers of stiffeners and different types of boundary conditions
    \item Run parametric analysis, vary input parameters
    \item Compare one analysis case with the results in Abaqus
    \item Present and discuss obtained results
\end{itemize}

Decision variables are the positions of stiffeners. The plate is rectangular. Adopt material data.
\vfill

\end{frame}
% ==================================================================


% ==================================================================
% {\frame{\frametitle{Inhalt}\tableofcontents}}
%\titleframe

\end{document}
% ==================================================================
